\section{Results}
\label{sec:results}

\subsection{Validation Overview}

We validate IBD detection and local ancestry inference against gold-standard VCF-based tools on chromosomes 10, 11, 12, and 20 of the HPRC dataset. Validation uses 223 HPRC individuals with matched reference panels, matched genetic maps, and the same samples for both tools, isolating methodological differences from compositional effects. IBD is compared against hap-ibd~\cite{browning2020hap_ibd}; local ancestry is compared against RFMix v2~\cite{maples2013rfmix}. Detailed parameter sweeps, per-pair breakdowns, and supplementary analyses are in Supplementary Section~K.

\subsection{IBD Detection}
\label{sec:results_ibd_validation}

\subsubsection{Baseline Concordance (chr20)}

On chr20 (44 sample pairs, 51 hap-ibd segments), the best HMM configuration (identity floor = 0.9, minimum segment length 2~Mb, 10 Baum-Welch iterations) achieves F1 = \textbf{0.283} (precision 0.317, recall 0.270). Scaling to all 2,850 HPRC pairs confirms the pattern: pair-level recall is 68.8\% (11/16 hap-ibd pairs detected) with precision of 4.3\% at the default threshold, improving to segment F1 = 0.432 at a 3.0~Mb minimum length (Supplementary Section~K.6).

The moderate concordance reflects a fundamental signal-to-noise limitation: the identity gap between IBD and non-IBD pairs ($\Delta\mu \approx 0.003$) is dwarfed by alignment-driven variance ($\sigma \approx 0.29$), yielding per-window SNR $\approx$ 0.01---approximately 55$\times$ worse than VCF-based methods (Supplementary Section~K.3). Eight independent algorithmic improvements (distance-aware transitions, genetic-map transitions, logit-transform emissions, background filtering, z-score normalization, population-adaptive transitions, coverage-ratio auxiliary emission, region exclusion) all converge to the same performance ceiling, establishing that \textbf{F1 $\approx$ 0.28 represents a measurement limitation, not a modeling limitation} (Supplementary Section~K.4).

\subsubsection{Multi-Chromosome Ranking-Based Evaluation}
\label{sec:multi_chr_ground_truth}

To test whether pangenome identity discriminates IBD from non-IBD pairs across diverse genomic contexts, we evaluated all 11 high-confidence hap-ibd pairs (LOD $\geq$ 3) across six regions spanning chr10, chr11, and chr12 (Table~\ref{tab:ranking_summary}).

\begin{table}[ht]
\centering
\caption{Multi-chromosome IBD ranking evaluation. For each region, all haplotype pairs are ranked by summary statistic of pangenome identity. ``Top 10\%'' indicates whether the IBD pair ranks in the top decile.}
\label{tab:ranking_summary}
\begin{tabular}{lcc}
\hline
\textbf{Summary statistic} & \textbf{IBD pairs in top 10\%} & \textbf{Ranked \#1} \\
\hline
Mean & 10/11 (91\%) & 6/11 \\
Trimmed mean (5\%) & \textbf{11/11 (100\%)} & 6/11 \\
Median & \textbf{11/11 (100\%)} & 6/11 \\
\hline
\end{tabular}
\end{table}

With robust statistics (trimmed mean or median), \textbf{all 11 IBD pairs (100\%) rank in the top 10\%} of pairs by pangenome identity, with \textbf{6 ranking \#1} in their region. The single outlier under the arithmetic mean (HG01175 vs.\ HG01192, chr10, rank 103/120) is caused by 5 alignment-artifact windows; robust statistics resolve it to rank 10/120 (top 8.3\%). At region-specific identity thresholds, \textbf{5 of 6 regions achieve perfect precision and recall} (Supplementary Section~K.11).

The identity gap between IBD and non-IBD pairs is remarkably consistent across chromosomes: $\Delta\mu = 0.00315$ (chr10), 0.00289 (chr11), 0.00326 (chr12)---all within $0.003 \pm 0.0004$, consistent with the chr20 finding. Segment length strongly discriminates true positives ($>$3.8~Mb) from false positives ($<$2.5~Mb).

\subsubsection{Distance-Aware Transitions}

On multi-chromosome HPRC validation, enabling distance-aware transitions with identity floor 0.998 improves mean pair-level F1 from 0.380 to \textbf{0.404} ($+$6.3\%), with per-chromosome results of chr11 F1 = 0.624 ($+$7.6\%), chr12 F1 = 0.440 ($+$6.8\%), and chr10 F1 = 0.148 (unchanged). This is now the recommended default for IBD detection.

\subsection{Local Ancestry Inference}
\label{sec:results_ancestry_validation}

\subsubsection{2-Way Baseline (chr20)}

On chr20 with a 2-way EUR/AFR model (HG00733, 30 EUR + 49 AFR reference panel), the \texttt{max} emission model achieves \textbf{93.2\%} concordance with RFMix (EUR F1 = 0.965, AFR F1 = 0.158). Multi-sample validation across 15 AMR individuals yields mean concordance \textbf{93.4\%} (range 90.2--95.5\%), confirming stability. The per-window signal-to-noise ratio for ancestry discrimination is SNR $\approx$ 0.09: in 75\% of true AFR windows, EUR references have \emph{higher} similarity than AFR references, making correct assignment statistically impossible at single-window resolution. The HMM achieves its concordance through temporal smoothing (transition-derived state persistence), not per-window signal detection (Supplementary Section~K.9).

\subsubsection{3-Way Concordance (chr10, chr11, chr12)}
\label{sec:3way_ancestry}

We extended to a 3-way EUR/AFR/AMR model on chr10, chr11, and chr12, modeling the Native American component present in all AMR query samples. Nine AMR query samples (18 haplotypes) were evaluated across six 1~Mb regions using a matched panel of 31 HPRC individuals (8 EUR, 8 AFR, 15 AMR). A sweep of 13 emission model configurations identified the top-3 averaging model as optimal (Table~\ref{tab:3way_results}).

\begin{table}[ht]
\centering
\caption{3-way ancestry concordance with RFMix (100 haplotype-region comparisons with available data, 6 regions across chr10/11/12). Top-3 emission model.}
\label{tab:3way_results}
\begin{tabular}{lccc}
\hline
\textbf{Ancestry} & \textbf{Precision} & \textbf{Recall} & \textbf{F1} \\
\hline
AMR & 95.5\% & 71.1\% & \textbf{0.815} \\
EUR & 29.2\% & 54.9\% & 0.381 \\
AFR & 30.8\% & 100.0\% & 0.471 \\
\hline
\multicolumn{3}{l}{Overall concordance} & \textbf{76.45\%} \\
\multicolumn{3}{l}{Per-chromosome range} & 71.9--80.7\% \\
\multicolumn{3}{l}{Best region (chr11:25--26~Mb)} & 88.9\% \\
\hline
\end{tabular}
\end{table}

The dominant ancestry (AMR) is detected with F1 = 0.815; minority ancestries are weaker, reflecting the signal-to-noise limitation of 62 pangenome reference haplotypes versus RFMix's 1,486 VCF haplotypes. The dominant error pattern is AMR windows misclassified as EUR or AFR (Supplementary Section~K.10).

\subsection{Simulation Validation}
\label{sec:sim_ancestry}

To evaluate under controlled conditions, we constructed a simulation framework using \texttt{msprime} to generate known ancestry tracts stitched onto real HPRC assemblies (\texttt{pangenome\_sim}), preserving realistic pangenome identity characteristics.

\subsubsection{Ancestry}

\begin{table}[ht]
\centering
\caption{Simulated ancestry concordance: \texttt{impop\_k} vs.\ RFMix on the same 46-haplotype panel (chr12, 5 individuals, pairwise contrast emissions).}
\label{tab:sim_ancestry}
\begin{tabular}{lccc}
\hline
\textbf{Metric} & \textbf{\texttt{impop\_k}} & \textbf{RFMix (46 hap)} & \textbf{RFMix (26 hap)} \\
\hline
Concordance & \textbf{97.95\%} & 95.9\% & 96.1\% \\
AFR F1 & 0.982 & 0.963 & 0.965 \\
EUR F1 & 0.977 & 0.954 & 0.956 \\
Gap vs.\ RFMix & \multicolumn{2}{c}{\textbf{+2.05 pp}} & --- \\
\hline
\end{tabular}
\end{table}

With pairwise contrast emissions, \texttt{impop\_k} achieves \textbf{97.95\%} concordance with ground truth, \emph{exceeding} RFMix (95.9\%) by \textbf{2.05\,pp} on the same panel. The pairwise contrast model---treating each population pair independently with per-pair adaptive temperature---is the key driver of improvement over the baseline softmax (95.1\% $\to$ 97.95\%).

\subsubsection{IBD}
\label{sec:sim_ibd}

On a simulated 50-individual pedigree (chr12), \texttt{impop\_k} achieves pair-level F1 = \textbf{0.514}, competitive with hap-ibd's 0.489 on the same data. Baum-Welch training is essential: without EM iterations, IBD detection fails entirely; with 10 iterations, the HMM converges to functional emission parameters. Panel size is critical: 30 reference haplotypes yield 0\% IBD detection, while 50 haplotypes achieve 51.4\%.

\subsection{Platinum Pedigree Validation (CEPH 1463)}
\label{sec:platinum_pedigree}

The CEPH 1463 (Platinum) family provides a real-data validation with known family structure: 5 of the 11 offspring have HPRC v2 assemblies, along with all 4 grandparents. We performed 4-state founder painting---assigning each genomic window in each offspring to one of the 4 grandparental origins based on pangenome IBS similarity.

The 4-state model achieves \textbf{99.53\%} mean concordance across all 5 offspring (range 97.5--99.5\%), making this the most convincing validation result. This is a real family with no simulation artifacts, demonstrating that pangenome assembly alignments alone can near-perfectly reconstruct grandparental inheritance patterns. The 8-state model (distinguishing maternal/paternal grandparental haplotypes) drops to 70--85\%, reflecting the inherent difficulty of resolving which haplotype within a grandparent was transmitted.

\subsection{Computational Performance}
\label{sec:computational_performance}

We benchmarked two IBS computation modes: the original per-window \texttt{impg} pipeline and PAF-direct mode (\texttt{ibs-from-paf}).

For a single-chromosome analysis (chr12, 40 reference haplotypes), \texttt{ibs-from-paf} completes in \textbf{46.3\,s} vs.\ 368\,s (7.9$\times$ speedup), with peak memory of \textbf{101\,MB} vs.\ 2,900\,MB (28.7$\times$ reduction). Scanning all 22 autosomes, \texttt{ibs-from-paf} processes the complete genome in \textbf{195\,s} (3.3\,min), yielding a \textbf{1,090$\times$} speedup over the \texttt{impg} pipeline ($\sim$59\,h extrapolated). Memory remains constant at 101\,MB regardless of genome coverage.

The complete VCF-based IBD/ancestry pipeline typically requires 5--7 hours per sample. The pangenome pipeline achieves comparable or superior accuracy in 3.3\,minutes of marginal compute time per sample pair, given pre-existing alignments---a fundamentally different point in the accuracy--speed trade-off space.
